\newpage{}
\section{Distribution Field and Observer Operators}
\label{sec:distribution-field-and-observer-operators}

\subsection{Edge Distribution Field}
\label{subsec:edge-distribution-field}

The MDN output for one sample is not just a point estimate.  It is a predictive
distribution attached to an edge, a channel, and a forecast horizon.  The
collection of these predictive laws over the active edge universe is the
\emph{edge distribution field}.

For a fixed forecast step and active edge universe, write
\[
  \mathfrak{D}^{c\,i}
  =
  \{\mathcal{D}_{e}^{c\,i} \in \mathcal{P}(\R^{D_f}) : e \in \mathcal{E}\}
  \in
  \left(\mathcal{P}(\R^{D_f})\right)^L.
\]
If the model is evaluated repeatedly over time, then for $t \in \mathcal{T}$,
$t \mapsto \mathfrak{D}^{c\,i}(t)$ is a time-indexed field of predictive
distributions.  The field itself stores uncertainty.  It is not yet a single
action, score, or node state.

\subsection{Observer Operators}
\label{subsec:observer-operators}

An observer operator is any map that reads a predictive distribution and returns
a value that can be acted on or compared.  Formally,
\[
  d_{\mathcal{O}} \in \{1,D_f\},
  \qquad
  \mathcal{O} : \mathcal{P}(\R^{d_{\mathcal{O}}}) \rightarrow \mathcal{Y},
\]
where $d_{\mathcal{O}}=1$ denotes a scalar observer and
$d_{\mathcal{O}}=D_f$ denotes an observer on the full future target vector.
Here $\mathcal{Y}$ is the chosen output space.

For $\mathcal{D} \in \mathcal{P}(\R^{D_f})$ and
$\vect{f} \in \R^{D_f}$, important coordinatewise examples are:
\begin{align*}
  \mathcal{O}_{\E}[\mathcal{D}] &= \E_{\mathcal{D}}[\vect{f}] \in \R^{D_f} &&\text{(expected value)},\\
  \mathcal{O}_{\Var}[\mathcal{D}] &= \Var_{\mathcal{D}}[\vect{f}] \in \R^{D_f} &&\text{(coordinatewise variance or scale)},\\
  \mathcal{O}_{\Quant\,\tau}[\mathcal{D}] &= \Quant_{\tau}(\mathcal{D}) \in \R^{D_f} &&\text{(coordinatewise $\tau$-quantile)},\\
  \mathcal{O}_{\ES\,\tau}[\mathcal{D}] &= \ES_{\tau}(\mathcal{D}) \in \R^{D_f} &&\text{(coordinatewise expected shortfall)},\\
  \mathcal{O}_{\Risk}[\mathcal{D}] &\in \R^m &&\text{(generic risk observer)}.
\end{align*}
For the expected-shortfall observer we follow the coherent-risk treatment of
Acerbi and Tasche~\cite{acerbi2002coherence}.

In plain terms, observers do not create the field.  They read the field.  The
expected value attached to an instrument edge is therefore one observer among
several, not the whole predictive object.

\subsection{Pointwise Fields Derived from Observers}
\label{subsec:pointwise-fields-derived-from-observers}

Applying an observer pointwise over the edge distribution field gives a derived
field.  For example, the mean field is
\[
  \widehat{r}_{e}^{c\,i}
  :=
  \mathcal{O}_{\E}[\mathcal{D}_{e}^{(r)\,c\,i}]
  \in \R,
\]
and a simple variance-based confidence rule is
\[
  \gamma_{e}^{c\,i}
  :=
  \frac{1}{\mathcal{O}_{\Var}[\mathcal{D}_{e}^{(r)\,c\,i}] + \delta}
  \in \Rpos,
  \qquad
  \delta \in \Rpos.
\]
The same idea yields downside-risk fields, entropy fields, or policy-specific
confidence fields.

\subsection{Second-Order Observers and Normalization}
\label{subsec:second-order-observers-and-normalization}

If a joint second-order observer is maintained across a collection of edge or
node returns $R_1,\ldots,R_m \in \R$, define the covariance matrix by
\[
  \Sigma_{i\,j} = \Cov(R_i,R_j),
  \qquad
  \Sigma \in \R^{m \times m}.
\]
The diagonal terms $\Sigma_{i\,i}$ are variances; they count one dispersion value
per coordinate, not necessarily one unit per traded instrument.  The normalized
correlation matrix is
\[
  \rho_{i\,j} = \frac{\Sigma_{i\,j}}{\sqrt{\Sigma_{i\,i}\Sigma_{j\,j}}},
  \qquad
  \rho \in \R^{m \times m}.
\]
This normalization removes scale and keeps only the signed strength of linear
co-movement.
